% Conclusion - ICML style

We study trip-level route generation and identify corridor-level multi-modality as a structural challenge for end-to-end coordinate generators and candidate-set formulations.
CascadeTraj addresses this by decomposing route generation into a short-horizon waypoint decision stage and an A*-based execution stage, making corridor commitment explicit on the road graph.

Our diagnostic gates establish that (i) heuristic $K$-shortest candidates can miss ground-truth corridors, motivating learned commitments; and (ii) time-of-day and road-tier context carry measurable signal for corridor variability in multi-modal OD groups.
Under these gates, CascadeTraj achieves near-perfect route completion with improved corridor match over node-level autoregressive baselines.

\paragraph{Limitations and future work.}
The current waypoint predictor uses coarse $32 \times 32$ bins, limiting corridor resolution.
Finer binning, adaptive waypoint counts, or end-to-end waypoint learning could improve corridor fidelity.
Reachability-aware decoding---incorporating graph connectivity into the decision stage---would further reduce execution failures.
Scaling to more cities and full WorldTrace, and connecting corridor generation to continuous executors for realistic within-corridor trajectories, are natural next steps.
